% !TEX root = ../algorithms.tex

\section{Онлайн-алгоритмы. Задача кэширования. Помечающий алгоритм}

\begin{Definition}{}{}
	\emph{Онлайн-алгоритм} --- алгоритм, получающий входные данные по одному элементу и обязанный реагировать на каждый элемент (принимать необратимое решение) до получения следующего, не зная будущих элементов. Для сравнения, \emph{офлайн-алгоритм} знает всю последовательность заранее.
\end{Definition}

\subsection{Задача кэширования}

\begin{Definition}{}{}
	Кэш состоит из $k$ ячеек (<<страниц>>). Пользователь последовательно запрашивает страницы $S_1, S_2, \ldots \in \{1,\ldots,N\}$, где $N > k$. Если запрошенная страница $S_i$ находится в кэше --- это \emph{попадание} (hit, стоимость $0$). Иначе --- \emph{промах} (miss, стоимость $1$): нужно загрузить $S_i$ в кэш, при необходимости вытеснив одну из $k$ находящихся там страниц. Цель онлайн-алгоритма --- минимизировать число промахов.
\end{Definition}

\begin{Definition}{}{}
	\emph{Оптимальный алгоритм} $\mathrm{OPT}$ --- алгоритм, решающий ту же задачу, но знающий всю последовательность запросов заранее и совершающий минимально возможное число промахов.
\end{Definition}

\begin{Definition}{}{}
	Онлайн-алгоритм называется \emph{$c$-конкурентным} ($c$-competitive), если на любой последовательности запросов число его промахов не превышает $c \cdot \mathrm{OPT}$ (плюс, возможно, аддитивная константа, зависящая от начального содержимого кэша), где $\mathrm{OPT}$ --- число промахов оптимального алгоритма на той же последовательности.
\end{Definition}

\begin{Definition}{}{}
	Алгоритм $\mathrm{LRU}$ (Least Recently Used): при промахе вытесняет из кэша страницу, к которой дольше всего не было обращений.
\end{Definition}

\begin{Theorem}{}{}
	Алгоритм $\mathrm{LRU}$ является $k$-конкурентным. Константу $k$ улучшить нельзя: существует последовательность, на которой $\mathrm{LRU}$ совершает в $k$ раз больше промахов, чем $\mathrm{OPT}$.
\end{Theorem}

\begin{proof}[Доказательство нижней оценки]
	Рассмотрим кэш размера $k$ и циклическую последовательность запросов к $k+1$ различным страницам:
	\[
	1, 2, \ldots, k, k{+}1,\; 1, 2, \ldots, k, k{+}1,\; \ldots
	\]
	После первых $k$ запросов (заполнения кэша) $\mathrm{LRU}$ промахивается на каждом следующем шаге: запрашиваемая страница всегда оказывается той, что была вытеснена как наименее недавно использованная на предыдущем шаге. Алгоритм $\mathrm{OPT}$, зная будущее, на той же последовательности совершает лишь один промах за каждые $k$ запросов (вытесняя страницу, которая понадобится позже остальных). Отношение числа промахов равно $k$, то есть $\mathrm{LRU}$ не может быть $c$-конкурентным ни для какого $c < k$.
\end{proof}

\subsection{Помечающий алгоритм}

\begin{Definition}{}{}
	Алгоритм \textsc{Marking} поддерживает для каждой страницы в кэше метку $0$ или $1$. При запросе страницы $s$ выполняется следующая процедура.
\end{Definition}

\begin{algorithmic}
	\Procedure{Mark}{$s$}
	\If{$s$ находится в кэше}
	\State пометить страницу $s$ единицей и завершить обработку запроса
	\Else
	\If{все страницы в кэше помечены единицей}
	\State сбросить метки всех страниц в кэше в ноль \Comment{начало нового периода}
	\EndIf
	\State выбрать случайную страницу, помеченную нулём, и вытеснить её из кэша
	\State поместить $s$ в кэш и пометить единицей
	\EndIf
	\EndProcedure
\end{algorithmic}

\begin{Theorem}{}{}
	Пусть $H_k = \sum_{i=1}^{k} \dfrac{1}{i} \approx \ln k$. Для любой последовательности запросов $S_1, S_2, \ldots$ математическое ожидание числа вытеснений алгоритма \textsc{Marking} не более чем в $2H_k$ раз превосходит число вытеснений алгоритма $\mathrm{OPT}$ на той же последовательности:
	\[
	\mathbf{E}\bigl[\#\text{вытеснений}_{\textsc{Marking}}\bigr]
	\;\leq\; 2H_k \cdot \mathrm{OPT}.
	\]
\end{Theorem}

\subsubsection*{Анализ по периодам}

\begin{Definition}{}{}
	Разобьём последовательность запросов $S_1, S_2, \ldots$ слева направо на \emph{периоды}: $p$-й период --- это максимальный отрезок запросов $S_i, S_{i+1}, \ldots, S_{i+\ell}$ такой, что
	\[
	\bigl|\{S_i,\ldots,S_{i+\ell}\}\bigr| = k,
	\qquad
	\bigl|\{S_i,\ldots,S_{i+\ell+1}\}\bigr| = k+1,
	\]
	то есть очередной запрос $S_{i+\ell+1}$ добавляет $(k+1)$-ю различную страницу.
\end{Definition}

\begin{Example}{}{}
	Для $k = 2$ последовательность
	\[
	a,\,b,\;c,\,c,\,b,\;e,\,e,\,s,\;f,\,f
	\]
	разбивается на периоды $a,b \ \big|\ c,c,b \ \big|\ e,e,s \ \big|\ f,f$.
\end{Example}

\begin{Definition}{}{}
	Пусть $A_p$ --- множество страниц, запрашиваемых в $p$-м периоде, $|A_p| = k$. Положим
	\[
	m_p = |A_p \setminus A_{p-1}|
	\]
	--- число <<новых>> страниц $p$-го периода (не встречавшихся в $(p-1)$-м). Так как $|A_p| = |A_{p-1}| = k$, также $|A_{p-1} \setminus A_p| = m_p$ и $|A_p \cap A_{p-1}| = k - m_p$.
\end{Definition}

\begin{Remark}{}{}
	В конце каждого периода (и только в эти моменты) все страницы кэша помечены единицей, после чего происходит сброс меток в ноль. Поэтому в начале $p$-го периода кэш содержит в точности $k$ страниц множества $A_{p-1}$ --- и только их.
\end{Remark}

\subsubsection*{Нижняя оценка для $\mathrm{OPT}$}

\begin{Lemma}{}{}
	$\mathrm{OPT} \geq \dfrac{1}{2}\displaystyle\sum_{p} m_p$.
\end{Lemma}

\begin{proof}
	Объединение $A_{p-1} \cup A_p$ содержит $|A_{p-1} \cup A_p| = |A_{p-1}| + |A_p \setminus A_{p-1}| = k + m_p$ различных страниц. Поскольку кэш вмещает только $k$ страниц, на объединённом отрезке запросов $(p-1)$-го и $p$-го периодов любой алгоритм, в частности $\mathrm{OPT}$, совершает не менее $m_p$ промахов.
	
	\begin{figure}[H]
		\centering
		\begin{tikzpicture}[>=Stealth, font=\small, box/.style={draw,minimum width=0.8cm,minimum height=0.6cm}]
			\foreach \i/\x in {1/0,2/1,3/2,4/3,5/4,6/5}
			\node[box] (p\i) at (\x,0) {$\i$};
			
			\draw (p1.north) -- (0,0.5) -- (1,0.5) -- (p2.north);
			\node at (0.5,0.78) {$m_2$};
			\draw (p3.north) -- (2,0.5) -- (3,0.5) -- (p4.north);
			\node at (2.5,0.78) {$m_4$};
			\draw (p5.north) -- (4,0.5) -- (5,0.5) -- (p6.north);
			\node at (4.5,0.78) {$m_6$};
			
			\draw (p2.south) -- (1,-0.5) -- (2,-0.5) -- (p3.south);
			\node at (1.5,-0.78) {$m_3$};
			\draw (p4.south) -- (3,-0.5) -- (4,-0.5) -- (p5.south);
			\node at (3.5,-0.78) {$m_5$};
		\end{tikzpicture}
		\caption{Периоды (пронумерованы $1,\ldots,6$) и оценки промахов $\mathrm{OPT}$. Каждая пара соседних периодов $(p-1,p)$ вместе затрагивает $k+m_p$ различных страниц, поэтому даёт $\mathrm{OPT}$ не менее $m_p$ промахов. Каждый период входит в две пары, $(p-1,p)$ и $(p,p+1)$, так что суммирование учитывает каждый промах $\mathrm{OPT}$ не более двух раз.}
	\end{figure}
	
	Каждый период $p$ участвует в двух таких парах --- $(p-1,p)$ и $(p,p+1)$, поэтому суммирование оценок по всем $p$ учитывает каждый промах $\mathrm{OPT}$ не более двух раз, откуда $\sum_{p} m_p \leq 2\,\mathrm{OPT}$.
\end{proof}

\subsubsection*{Верхняя оценка для \textsc{Marking}}

\begin{Lemma}{}{}
	Рассмотрим $p$-й период в худшем для \textsc{Marking} порядке: первые $m_p$ запросов --- к страницам множества $A_p \setminus A_{p-1}$ (<<новые>>), а оставшиеся $k - m_p$ запросов --- к страницам $S^{(1)}, \ldots, S^{(k-m_p)}$, образующим множество $A_p \cap A_{p-1}$ (<<старые>>), в произвольном порядке. Тогда для каждого $j = 0, \ldots, k-m_p-1$ вероятность промаха на запросе $S^{(j+1)}$ не превышает $\dfrac{m_p}{k-j}$.
\end{Lemma}

\begin{proof}
	По предыдущему замечанию в начале $p$-го периода кэш состоит из всех $k$ страниц множества $A_{p-1}$, и все они не помечены. Каждый из первых $m_p$ запросов --- промах: алгоритм вытесняет случайную непомеченную страницу. Так как ни одна из старых страниц ещё не помечена, эти $m_p$ вытеснений выбирают (без повторений) равномерно случайное $m_p$-подмножество $E \subseteq A_{p-1}$.
	
	Зафиксируем $j$ и предположим, что $S^{(1)},\ldots,S^{(j)} \notin E$ (то есть к моменту своих запросов они не были вытеснены). По симметрии множество $E$ равномерно распределено среди оставшихся $k-j$ страниц $A_{p-1}$, поэтому
	\[
	\Pr\bigl(S^{(j+1)} \in E \mid S^{(1)},\ldots,S^{(j)} \notin E\bigr) = \frac{m_p}{k-j}.
	\]
	Если $S^{(j+1)} \in E$, то страница уже вытеснена из кэша к моменту своего запроса, и этот запрос --- промах. Следовательно, вероятность промаха на $S^{(j+1)}$ не превышает $\dfrac{m_p}{k-j}$.
\end{proof}

\begin{Lemma}{}{}
	$\mathbf{E}\bigl[\#\text{промахов в }p\text{-м периоде}\bigr] \leq m_p H_k$.
\end{Lemma}

\begin{proof}
	В худшем порядке из предыдущей леммы первые $m_p$ запросов периода --- гарантированные промахи. Для остальных запросов положим
	\[
	I_j =
	\begin{cases}
		1, & (j{+}1)\text{-й <<старый>> запрос --- промах},\\
		0, & \text{иначе},
	\end{cases}
	\qquad j = 0,\ldots,k-m_p-1.
	\]
	Тогда число промахов в $p$-м периоде равно $m_p + \sum_{j=0}^{k-m_p-1} I_j$, и по предыдущей лемме
	\begin{align*}
		\mathbf{E}\!\left[m_p + \sum_{j=0}^{k-m_p-1} I_j\right]
		&\leq m_p + \sum_{j=0}^{k-m_p-1} \frac{m_p}{k-j}
		= m_p + m_p \sum_{h=m_p+1}^{k} \frac{1}{h}
		= m_p\left(1 + \sum_{h=m_p+1}^{k} \frac{1}{h}\right).
	\end{align*}
	Так как $H_{m_p} = \sum_{i=1}^{m_p} \tfrac{1}{i} \geq 1$ при $m_p \geq 1$, выполняется $\sum_{h=m_p+1}^{k} \tfrac{1}{h} = H_k - H_{m_p} \leq H_k - 1$, откуда
	\[
	\mathbf{E}\bigl[\#\text{промахов в }p\text{-м периоде}\bigr] \leq m_p\bigl(1 + (H_k - 1)\bigr) = m_p H_k.
	\]
	При $m_p = 0$ оценка тривиальна.
\end{proof}

\subsubsection*{Завершение доказательства}

\begin{proof}[Доказательство теоремы]
	По линейности математического ожидания и предыдущей лемме,
	\[
	\mathbf{E}\bigl[\#\text{вытеснений}_{\textsc{Marking}}\bigr]
	= \sum_p \mathbf{E}\bigl[\#\text{промахов в }p\text{-м периоде}\bigr]
	\leq \sum_p m_p H_k
	= H_k \sum_p m_p.
	\]
	По лемме о нижней оценке $\sum_p m_p \leq 2\,\mathrm{OPT}$, поэтому
	\[
	\mathbf{E}\bigl[\#\text{вытеснений}_{\textsc{Marking}}\bigr] \leq H_k \cdot 2\,\mathrm{OPT} = 2H_k \cdot \mathrm{OPT}.
	\]
\end{proof}

\begin{Remark}{}{}
	Известно также, что любой вероятностный алгоритм кэширования не может быть лучше $H_k$-конкурентного относительно $\mathrm{OPT}$; таким образом, оценка $2H_k$ для \textsc{Marking} оптимальна с точностью до множителя $2$.
\end{Remark}